% !TeX root = ../papers/nc-p3.tex
\section{射影した楕円五次・有理五次の階数一曲線の排除}
\label{sec:rank-one-projected-quintics}
\begin{theorem}\label{thm:rank-one-projected-quintics}
定理\ref{thm:rank-one-total-five}の被約窓では、滑らかな整五次曲線は階数一軌跡に存在しない。
さらに、正規化が有理で二次曲面に載らない特異な整五次曲線も存在しない。
\end{theorem}
\begin{proof}[二つの表示に共通する切断]
正規化上の二つの超平面束が同型で、それぞれの座標四空間 $U,V$ に独立な共通切断 $s,t$ があり、
入力曲線上で二次評価が単射とする。$s^2,st,t^2$ の三つに対し二つの表示が同じ切断を与えるから、
反対根による二次座標同型もこの三つをそのまま対応させる。実際共通の二乗束の同一視の違いは定数倍だけで、
入力側の二次多項式は曲線への制限で一意である。従って各印で対応する二成分ずつの平方が比例し、
双一次式 $s_x t_y-t_x s_y$ が全二十印で消える。これは十六の混合積の独立性に反する。
五次有理正規化では六切断の中の二つの四空間が少なくとも二次元で交わる。また二つの束は必ず同型である。
これで二次曲面に載らない有理正規化の場合を除いた。

滑らかな有理五次が二次曲面に載れば、両表示とも滑らかな二次曲面上の $(1,4)$ 曲線である。
実際二次式の核次元は両側で等しい。錐の解消 $\mathbb F_2$ 上の真変換 $bE+5f$ は
滑らかな像と同型であり、頂点を通れば $5-2b=1$、避ければ $5-2b=0$。
可能なのは $b=2$ で種数 $(b-1)(5-b-1)=2$ のみである。
滑らかな二次曲面では二次数の和5と種数公式から結論する。
従って各座標四空間は二項式環の $S_1 W$ の形で、$W\subset S_4$ は基点のない鉛筆である。
その二乗 $S_2 W^2\subset S_{10}$ から $W$ を復元できる。
実際 $W=\langle f,g\rangle$ とすると $I=(f,g)^2$ の自由分解は
$0\to S(-12)^2\to S(-8)^3\to I\to0$ である。二生成正則列の二乗の関係は
$g f^2-f(fg)$ と $g(fg)-f g^2$ で尽きる。従って $S/I$ のsocleは次数10のみにあり
（最上のKoszulホモロジーを二変数分ずらす）、$(I_{10}:S_2)_8=I_8=W^2$。
この三切断による射は $W$ の鉛筆の後に平面二次Veroneseを施したもので、像の滑らかな二次曲線上の
次数一束を使えば鉛筆も同じになる。両表示の二乗空間が等しいので $U=V$、
二つの五次埋込みは射影同値である。唯一の二次式を用いる前定理の $xx^t+qA$ の議論がこれを排除する。
\end{proof}
\begin{proof}[楕円五次の余核と長さ三因子]
滑らかな楕円五次 $E$ は二次曲面に載らない。滑らかな二次曲面で可能な二次数は $(1,4),(2,3)$ で種数は0または2、
錐では上の計算で種数2だからである。$L'=L\eta$ で $\eta=1$ なら
五切断中の二四空間の交わりは三次元以上で最初の議論に反する。$\eta\ne1$ とする。
各二乗座標空間は十次元の $H^0(2L)$ 全体である。曲線外の階数一点を
$[\lambda]\in\mathbb P H^0(2L)^*$ とする。$\lambda$ による積対合は、
$H^0(L)$ と $H^0(L')$ のそれぞれで、指定四空間への制限の階数が1なので全階数は高々3である。

Serre双対性で $\lambda$ を拡張 $0\to L^{-1}\to F\to L\to0$ の類として読む。
$F$ と $F\eta$ の連結写像は上の二つの積対合であり、いずれも核次元が2以上である。
$F$ が半安定ならその非零切断は零点を持たず、行列式が自明なので自明束同士の拡張となる。
第二の独立な切断により分裂し $H^0(F\eta)=0$ となるので矛盾する。
よって正次数の飽和直線束 $A\subset F$ があり、商は $A^{-1}$、$h^0(F)=h^0(A)=\deg A\ge2$。
$A\to L$ は非零で、$A=L(-D)$ と書けば有効因子 $D$ の長さは3以下である。
この包含が拡張に持ち上がるから、双対性で $\lambda$ は $H^0(2L(-D))$ を消す。
従って二乗埋込みで $[\lambda]\in\langle D\rangle$ である。

$D$ をこれを満たす最小の部分因子に取り直す。$\lambda$ の $D$ 上の積対合は非退化である
（局所単項環のsocleを消せば因子長を一つ減らせる）。長さ2以下なら各四空間は埋込みを与え
$D$ 上で全射なので、制限階数はその長さと等しい。曲線外なので長さは3でなければならない。
同様に二つの四空間は $D$ 上でともに階数2で、これは二つの埋込みでの共通三割因子である。
一つの $D$ から曲線外の点は高々一つだけ生じる。その二乗の包は両側で三割直線の二次Veroneseの張る平面である。
そこで二つの滑らかな二次曲線が一致すれば別の階数一直線ができ総次数5を超え、
異なれば長さ4の交叉のうち3を $D$ が既に使っている。
\end{proof}
\begin{proof}[三割因子の二つの切断]
原点 $o\in E$ を固定し $E\times E$ の第二射影を $p$、第一射影を $p_1$、対角を $\Delta$ とする。
$\mathcal B=p_1^*\mathcal O(o)(\Delta)$、$\mathcal A=p_1^*L\otimes\mathcal B^{-1}$ と書く。
底を動くと $\mathcal A$ の繊維は全ての次数三線束を一度ずつ走る。
$U=\ker\alpha\subset H^0(L)$ に対し積の評価
$p_*\mathcal A\to(p_*\mathcal B)^*$ は全射である。各繊維で、もし非零の二次線束の切断 $b$ が
$\alpha(b H^0(\mathcal A_t))=0$ を満たせば $U$ が $b$ の長さ二零因子で全射にならず埋込みに反する。
従って核は一本 $K$ であり、各繊維のその切断の零因子が唯一の三割因子である。実際長さ三因子 $D$ 上での
$U$ の評価が二次元になる条件は、$H^0(L(-D))$ の二次元全体を $U$ が含むことである。
$L'$ について同じ $\mathcal A$ を用い、$\mathcal B'=\mathcal B\otimes p_1^*\eta$ を取る。

対角の制限完全列で $p_*\mathcal A$ は $H^0(L(-o))\otimes\mathcal O\to L(-o)$ の評価核なので次数 $-4$。
$p_*\mathcal B$ は $\mathcal O$ を部分束、$\mathcal O(o)$ を商とし、$p_*\mathcal B'$ ではその商が $\mathcal O(o)\eta$ である。
これは $\mathcal O(\Delta)|_\Delta$ が楕円曲線の自明な接束であるためである。よって $\deg K=-3$。
二つの三割条件の一致は $K\to (p_*\mathcal B')^*$ の零点である。この射が恒等零なら両評価核は同じ部分束で、
商束の同型が行列式 $\mathcal O(-o)$ と $\mathcal O(-o)\eta$ の違いに反する。
$K^{-1}(p_*\mathcal B')^*$ は次数2の直線束を部分、次数3の直線束を商に持つため、
この非零切断の零点は重複込みで高々三つである。各因子が高々一つの追加点を与えるので曲線外の印は三つ以下。
しかし $h^0(LL')=10$ と混合積の独立性から六つ以上必要である。この矛盾と前定理で全ての滑らかな場合を終える
（非平面五次の種数はCastelnuovo上界2以下）。
\end{proof}
